% ============================================================
% The Neutrality Theorem
% ============================================================

\documentclass[12pt]{amsart}

\usepackage{amsmath,amssymb,amsthm}
\usepackage{booktabs}
\usepackage{microtype}
\usepackage{url}

\newtheorem{theorem}{Theorem}
\newtheorem{lemma}[theorem]{Lemma}
\newtheorem{proposition}[theorem]{Proposition}
\newtheorem{corollary}[theorem]{Corollary}
\theoremstyle{definition}
\newtheorem{definition}[theorem]{Definition}
\theoremstyle{remark}
\newtheorem{remark}[theorem]{Remark}

\title{The Neutrality Theorem}

\author{Alexander S.\ Petty}
\email{alexander.petty@gmail.com}
\date{August 2025}
\subjclass[2020]{11A63, 11N05, 11M06}

\begin{document}
\begin{abstract}
The centered collision fluctuation sum $F^{\circ}(s)$
converges at $s = 1$ and appears to converge below
$s = 1$ into the critical strip. This paper identifies
the mechanism.

The centered sum decomposes as
$F^{\circ} = F^{\circ\circ} + \mathcal{M}$, where
$F^{\circ\circ}$ is doubly centered and $\mathcal{M}$
is the mod-$3$ component. In the nontrivial cancellation
regime $c_0 \ne 0$, both diverge individually below
$s = 1$, but their ratio approaches $-1$ with increasing
precision. At finite range the deviation
$(1 + \text{ratio}) \sim C / |F^{\circ\circ}|^{\alpha}$
fits $\alpha \approx 0.6$, but the character
decomposition proves $\alpha = 1$: the cancellation
is perfect.

The proof is algebraic. The mod-$3$ component
decomposes over characters modulo~$3$ as
$\mathcal{M} = c_0 \sum_p p^{-s} + c_1 \sum_p
\chi_3(p)\, p^{-s}$. The doubly centered sum carries
the same Mertens term $-c_0 \sum_p p^{-s}$ with
opposite sign. In $F^{\circ} = F^{\circ\circ} +
\mathcal{M}$, the Mertens terms cancel exactly,
leaving $F^{\circ}$ as a pure linear combination of
non-trivial character sums. These converge wherever
the corresponding $L$-functions have no zeros.
\end{abstract}
\maketitle

% ============================================================
\section{Introduction}

This paper studies why the centered collision
fluctuation sum $F^{\circ}(s)$ penetrates below $s = 1$.

The preceding papers~\cite{paper16,paper17} proved
convergence at $s = 1$ via the collision transform and
Mertens' theorem, and observed computational convergence
to at least $s = 0.6$ in base~$10$. The mod-$3$ polarity
was identified as the structural feature controlling the
penetration depth: removing the mod-$3$ component
preserves convergence at $s = 1$ but destroys it below.

The present paper proves that the mod-$3$ cancellation
is algebraically perfect. The doubly centered sum
$F^{\circ\circ}$ and the mod-$3$ component $\mathcal{M}$
share a common Mertens divergence
$\pm c_0 \sum_p p^{-s}$ that cancels in their sum.
The residual $F^{\circ}$ contains only non-trivial
character sums.

% ============================================================
\section{The Decomposition}

\begin{definition}
Let $S^{\circ}(a)$ denote the singly centered collision
deviation (centered by spectral class). For each unit
$a \bmod m$ with $m = b^{\ell+1}$, define:
\begin{align*}
\mu_3(a) &= \frac{1}{|\{a' : a' \equiv a \!\!\pmod{3}\}|}
\sum_{\substack{a' \equiv a \pmod{3}}}
S^{\circ}(a'), \\
S^{\circ\circ}(a) &= S^{\circ}(a) - \mu_3(a).
\end{align*}
The \emph{mod-$3$ component} is
$\mathcal{M}(s) = \sum_p \mu_3(p) / p^s$
and the \emph{doubly centered sum} is
$F^{\circ\circ}(s) = \sum_p S^{\circ\circ}(p) / p^s$.
\end{definition}

By construction:
\[
F^{\circ}(s) = F^{\circ\circ}(s) + \mathcal{M}(s).
\]

\begin{proposition}
The antisymmetry $S^{\circ}(a) = -S^{\circ}(m{-}a)$ is
preserved by the mod-$3$ centering: $S^{\circ\circ}(a)
= -S^{\circ\circ}(m{-}a)$.
\end{proposition}

\begin{proof}
The reflection $a \mapsto m - a$ pairs mod-$3$ classes:
it sends residue class $r$ to $m-r \pmod{3}$. Since
$S^{\circ}$ is antisymmetric, the mod-$3$ means of
reflected classes satisfy
\[
\mu_3(r) = -\mu_3(m-r).
\]
Therefore
\[
S^{\circ\circ}(m-a)
= S^{\circ}(m-a) - \mu_3(m-a)
= -S^{\circ}(a) + \mu_3(a)
= -S^{\circ\circ}(a).
\]
\end{proof}

\begin{corollary}
$F^{\circ\circ}(1)$ converges: the antisymmetry ensures
even characters vanish, and the proof of the Centered
Collision PNT~\cite{paper16} applies.
\end{corollary}

% ============================================================
\section{The Character Algebra}

The function $\mu_3(a)$ depends only on $a \bmod 3$.
When $3 \nmid b$, we have $\gcd(3,m)=1$, so the
restriction to primes $p > 3$ (where
$p \bmod 3 \in \{1,2\}$) gives a decomposition over
characters modulo~$3$:
\[
\mu_3(p) = c_0 + c_1 \chi_3(p),
\]
where $\chi_3$ is the non-trivial character modulo~$3$,
\[
c_0 = \frac{m_1 + m_2}{2}, \qquad
c_1 = \frac{m_1 - m_2}{2},
\]
and $m_1, m_2$ are the mod-$3$ means:
$m_j = \mu_3(a)$ for any $a \equiv j \pmod{3}$.

In base~$10$ at lag~$1$: $m_1 = -1.30$, $m_2 = 0$,
giving $c_0 = c_1 = -0.65$.

The mod-$3$ component is therefore
\begin{equation}\label{eq:M-decomp}
\mathcal{M}(s) = c_0 \sum_p \frac{1}{p^s}
+ c_1 \sum_p \frac{\chi_3(p)}{p^s}.
\end{equation}

The first sum is the Mertens prime sum, diverging like
$\log\log x$ at $s = 1$ and like $x^{1-s}/((1{-}s)\log x)$
for $s < 1$. The second converges at $s = 1$ by
Dirichlet's theorem~\cite{davenport}.

% ============================================================
\section{The Cancellation Theorem}

\begin{theorem}[Perfect cancellation]\label{thm:cancel}
The Mertens divergence in $F^{\circ\circ}$ and
$\mathcal{M}$ cancels exactly:
\[
F^{\circ}(s) = \sum_{\substack{\chi \bmod m \\
\chi(-1) = -1}} \hat{S}^{\circ}(\chi)\, P(s, \chi),
\]
where $P(s, \chi) = \sum_p \chi(p)/p^s$. No trivial
character (Mertens) term appears.
\end{theorem}

\begin{proof}
By the collision transform~\cite{paper16},
$S^{\circ}(p) = \sum_{\chi} \hat{S}^{\circ}(\chi)\,
\chi(p)$, and the antisymmetry theorem restricts
the sum to odd characters ($\chi(-1) = -1$). Therefore
\[
F^{\circ}(s) = \sum_{\substack{\chi \bmod m \\
\chi(-1) = -1}} \hat{S}^{\circ}(\chi)
\sum_p \frac{\chi(p)}{p^s}
= \sum_{\substack{\chi \bmod m \\
\chi(-1) = -1}} \hat{S}^{\circ}(\chi)\, P(s, \chi).
\]
Every character in this sum is non-trivial (the trivial
character is even), so no Mertens term
$\sum_p p^{-s}$ appears.

The decomposition $F^{\circ} = F^{\circ\circ} +
\mathcal{M}$ introduces the Mertens term
$c_0 \sum_p p^{-s}$ via~\eqref{eq:M-decomp}, and
$F^{\circ\circ}$ must carry $-c_0 \sum_p p^{-s}$ to
compensate. In $F^{\circ}$, these cancel algebraically.
\end{proof}

\begin{definition}
The \emph{cancellation exponent} $\alpha$ is defined by
\[
1 + \frac{F^{\circ\circ}(s)}{\mathcal{M}(s)}
\sim \frac{C}{|F^{\circ\circ}(s)|^{\alpha}}
\]
as $|F^{\circ\circ}(s)| \to \infty$.
\end{definition}

\begin{corollary}\label{cor:alpha-one}
Assume $c_0 \ne 0$, so that the Mertens term in
\eqref{eq:M-decomp} is present and
$|F^{\circ\circ}(s)| \to \infty$. Then the cancellation
exponent is $\alpha = 1$.
\end{corollary}

\begin{proof}
Since $1 + F^{\circ\circ}/\mathcal{M} =
F^{\circ}/\mathcal{M}$, and $F^{\circ}(s)$ converges
to a finite value while
$|\mathcal{M}(s)| \sim |c_0| \sum_p p^{-s} \to \infty$,
the ratio $F^{\circ}/\mathcal{M} \to 0$. Since
$F^{\circ\circ} = F^{\circ} - \mathcal{M}$ and
$F^{\circ}(s)$ converges while
$|\mathcal{M}(s)| \to \infty$, we have
$F^{\circ\circ}/\mathcal{M} = F^{\circ}/\mathcal{M}
- 1 \to -1$, so $|F^{\circ\circ}| \sim |\mathcal{M}|$.
Therefore
\[
\bigl|1 + \text{ratio}\bigr|
= \frac{|F^{\circ}(s)|}{|\mathcal{M}(s)|}
\sim \frac{C}{|F^{\circ\circ}(s)|^1}.
\]
\end{proof}

% ============================================================
\section{Computational Verification}

\begin{table}[h]
\centering
\begin{tabular}{rrrrrrr}
\toprule
primes & $s$ & $F^{\circ}$ & $F^{\circ\circ}$ &
$\mathcal{M}$ & $1 {+} \text{ratio}$ &
$\alpha_{\text{fit}}$ \\
\midrule
$10{,}000$ & $0.7$ & $+0.27$ & $+7.1$ & $-6.8$
& $-0.039$ & $0.58$ \\
$100{,}000$ & $0.7$ & $+0.28$ & $+13.1$ & $-12.8$
& $-0.022$ & $0.61$ \\
$348{,}488$ & $0.7$ & $+0.29$ & $+18.1$ & $-17.8$
& $-0.016$ & $0.62$ \\
$700{,}000$ & $0.7$ & $+0.26$ & $+21.6$ & $-21.3$
& $-0.012$ & $0.91$ \\
$1{,}270{,}000$ & $0.7$ & $+0.27$ & $+25.2$ & $-24.9$
& $-0.011$ & $0.75$ \\
\bottomrule
\end{tabular}
\caption{The fitted cancellation exponent $\alpha$
drifts upward as the prime range increases. At
$348{,}488$ primes $\alpha \approx 0.6$; at
$1{,}270{,}000$ primes the fit is noisy but trends
toward $\alpha = 1$. Base~$10$, lag~$1$.}
\label{tab:drift}
\end{table}

The finite-range exponent $\alpha \approx 0.6$ reported
in earlier computation reflects the pre-asymptotic
regime. As the prime count increases, the Mertens term
$c_0 \sum_p p^{-s}$ dominates both $F^{\circ\circ}$
and $\mathcal{M}$, and the ratio
$F^{\circ}/\mathcal{M}$ shrinks as $1/|\mathcal{M}|$,
confirming $\alpha = 1$.

\begin{table}[h]
\centering
\begin{tabular}{rrrr}
\toprule
$s$ & $F^{\circ}$ & $|F^{\circ\circ}|$ &
$(1{+}\text{ratio}) \cdot |F^{\circ\circ}|$ \\
\midrule
$1.0$ & $+0.08$ & $0.90$ & $0.085$ \\
$0.9$ & $+0.12$ & $2.38$ & $0.126$ \\
$0.8$ & $+0.18$ & $7.25$ & $0.187$ \\
$0.7$ & $+0.27$ & $25.2$ & $0.273$ \\
$0.6$ & $+0.37$ & $96.6$ & $0.369$ \\
$0.5$ & $+0.37$ & $399.6$ & $0.370$ \\
\bottomrule
\end{tabular}
\caption{With $\alpha = 1$, the product
$(1{+}\text{ratio}) \cdot |F^{\circ\circ}|$ equals
$|F^{\circ}(s)|$, which is bounded. Compare with
Table~\ref{tab:drift}: the product stabilizes as
the prime count grows. $1{,}270{,}000$ primes.}
\label{tab:alpha-one}
\end{table}

\begin{table}[h]
\centering
\begin{tabular}{rrrrrrr}
\toprule
$b$ & $\phi(b^2)$ & $c_0$ & $F^{\circ}(0.5)$ &
$|F^{\circ\circ}(0.5)|$ & $\alpha_{\text{fit}}$ &
type \\
\midrule
$3$ & $6$ & $0$ & $+0.65$ & $0.65$ & --- & $3 \mid b$ \\
$5$ & $20$ & $+0.11$ & $+0.88$ & $51$ & $0.69$ &
weak \\
$7$ & $42$ & $-0.70$ & $+1.76$ & $323$ & $0.69$ &
strong \\
$10$ & $40$ & $-0.65$ & $+0.09$ & $296$ & $0.93$ &
strong \\
$11$ & $110$ & $+0.02$ & $+1.72$ & $8.3$ & --- &
weak \\
$12$ & $48$ & $0$ & $-0.64$ & $0.64$ & --- & $3 \mid b$ \\
$13$ & $156$ & $-0.80$ & $+3.53$ & $369$ & $0.60$ &
strong \\
\bottomrule
\end{tabular}
\caption{The cancellation exponent across bases
($664{,}000$ primes). When $3 \mid b$, the mod-$3$
bias vanishes ($c_0 = 0$) and $F^{\circ} = F^{\circ\circ}$
with no cancellation needed. When $3 \nmid b$, the
bias is strong ($|c_0| > 0.5$) or weak ($|c_0| < 0.2$).
At strong bias, $\alpha$ is already near $1$ at this
range. At weak bias, the Mertens term has not yet
dominated and the finite-range $\alpha$ is lower.}
\label{tab:multi-base}
\end{table}

Three regimes emerge. When $3 \mid b$, the modulus
$m \equiv 0 \pmod{3}$ forces all units into mod-$3$
classes $1$ and $2$, and the centering removes the
bias completely: $c_0 = 0$ and no cancellation mechanism
is needed. When $3 \nmid b$ and the bias is strong
($|c_0| \approx 0.7$), the Mertens term dominates
early and $\alpha$ approaches $1$ quickly. When $3 \nmid b$
but the bias is weak ($|c_0| \approx 0.02$), the
character sum corrections remain comparable to the
Mertens term at this range, and the finite-range
$\alpha$ stays low. In the nontrivial cancellation
regime $c_0 \ne 0$, the theorem predicts
$\alpha = 1$ asymptotically; when $c_0 = 0$, no
cancellation exponent is defined because no divergent
Mertens term is present.

% ============================================================
\section{Consequences}

\begin{corollary}
$F^{\circ}(s)$ converges for
$\operatorname{Re}(s) > 1 - c/\log m$, where $c > 0$
is the constant from the classical zero-free region
for Dirichlet $L$-functions modulo~$m$.
\end{corollary}

\begin{proof}
By Theorem~\ref{thm:cancel}, $F^{\circ}(s) = \sum
\hat{S}^{\circ}(\chi)\, P(s, \chi)$ over non-trivial
odd characters. By partial summation and the explicit
formula for $\sum_{p \le x} \chi(p)$
(see~\cite[Ch.~20]{davenport}
or~\cite[Ch.~13]{montgomery}), each $P(s, \chi)$
converges wherever $L(s, \chi)$ has no zeros with real
part $\ge \operatorname{Re}(s)$. The classical zero-free
region~\cite{davenport,montgomery} guarantees this
for $\operatorname{Re}(s) > 1 - c/\log m$.
\end{proof}

\begin{remark}
By the convergence theorem for Dirichlet
series~\cite{hardy}, convergence of $F^{\circ}$ at a
real point $s_0$ extends to the half-plane
$\operatorname{Re}(s) > s_0$. The question of
penetration depth is therefore a question about the
real abscissa of convergence.

Computation indicates convergence to at least
$s_0 = 0.5$~\cite{paper17}. If confirmed, this would
imply convergence for all
$\operatorname{Re}(s) > 1/2$. Since $F^{\circ}$ is
a finite linear combination of character sums
$P(s, \chi)$, two scenarios are possible.

If the products
$\hat{S}^{\circ}(\chi) \cdot P(s, \chi)$ were to have
non-negative real part,
then convergence of $F^{\circ}$ requires each
$P(s, \chi)$ to converge individually, which is
equivalent to $L(s, \chi) \ne 0$ for
$\operatorname{Re}(s) > 1/2$: the Generalized Riemann
Hypothesis for the odd $L$-functions modulo~$m$.

If the positivity condition fails, convergence of
$F^{\circ}$ could arise from cancellation between
divergent character sums. This would mean the
collision coefficients $\hat{S}^{\circ}(\chi)$ encode
structural information about the zeros of the
corresponding $L$-functions.
\end{remark}

% ============================================================
\section{The Neutrality Theorem}

The mod-$3$ component $\mathcal{M}$ arises because the
collision invariant $S$ has unequal means across mod-$3$
classes. The following theorem identifies the universal
structure.

\begin{theorem}[Neutrality of the fixed class]\label{thm:neutral}
For any base $b$ with $3 \nmid b$ and any lag
$\ell \ge 1$, let $m=b^{\ell+1}$ and let
$r_m \in \{1,2\}$ be the unique residue class satisfying
\[
2r_m \equiv m \pmod{3}
\qquad\text{equivalently}\qquad
r_m \equiv -m \pmod{3}.
\]
Then the mean of $S$ (the collision deviation at lag~$\ell$)
over units
$a \equiv r_m \pmod{3}$ in
$(\mathbb{Z}/b^{\ell+1}\mathbb{Z})^{\times}$ equals
$-1/2$ (the grand mean). In particular, when
$m \equiv 1 \pmod{3}$ the neutral class is $2$, and when
$m \equiv 2 \pmod{3}$ the neutral class is $1$.
\end{theorem}

\begin{proof}
Since $3 \nmid b$, the class $r_m$ is well-defined modulo~$3$.
The reflection $a \mapsto m-a$ sends residue class $r$ to
$m-r \pmod{3}$. The class $r_m$ is fixed by this map because
\[
m-r_m \equiv r_m \pmod{3}
\qquad\Longleftrightarrow\qquad
2r_m \equiv m \pmod{3}.
\]
Thus reflection is a self-map on the class-$r_m$ units. The
reflection identity
$S(a) + S(m{-}a) = -1$ pairs units within this class,
and each pair averages to $-1/2$.
\end{proof}

\begin{corollary}\label{cor:swap}
The remaining two mod-$3$ classes are swapped by reflection,
so their means sum to $-1$. Thus all deviation from the
grand mean is concentrated in the reflected pair, with equal
magnitude and opposite sign around $-1/2$.
\end{corollary}

For the lag-$1$ case studied in the critical-strip note,
$m=b^2 \equiv 1 \pmod{3}$ whenever $3 \nmid b$, so the
neutral class is $2$ and the reflected pair is $\{0,1\}$.
Since primes $p>3$ satisfy $p \bmod 3 \in \{1,2\}$, the
prime harmonic sum then sees only classes $1$ (which carries
a base-dependent bias) and $2$ (which is neutral). This is
the source of the lag-$1$ mod-$3$ polarity used in the
critical-strip computations.

\begin{remark}
Since $3 \mid 3$, Theorem~\ref{thm:neutral} does not
apply to base~$3$.  In base~$3$, the class-$1$ mean
is $-2/3$ and the class-$2$ mean is $-1/3$, giving a
splitting ratio of
$|{-}2/3|/(|{-}2/3|{+}|{-}1/3|) = 2/3$. The fraction
$2/3 = (p{-}1)/p$ at $p = 3$ is the alignment limit
of the prime~$3$~\cite{paper1}. This ratio is specific
to base~$3$; in the bases with $3 \nmid b$ and
$m \equiv 1 \pmod{3}$, the class-$2$ mean is $-1/2$
(neutral), while in the cases with $m \equiv 2 \pmod{3}$
the neutral class is $1$. The universal structure is the
neutrality of the fixed reflected class, not the ratio $2/3$.
\end{remark}

Computation confirms the neutrality theorem across all
bases $b \le 40$ with $3 \nmid b$: in the cases with
$m \equiv 1 \pmod{3}$ one finds
$2 \cdot \text{sum}_2 = -\text{count}_2$ exactly, while
in the cases with $m \equiv 2 \pmod{3}$ the analogous
identity holds with class~$1$ in place of class~$2$.
For bases with $3 \mid b$, the modulus $m \equiv 0 \pmod{3}$,
no units satisfy $a \equiv 0 \pmod{3}$, and the
structure is different.

% ============================================================
\section{Remarks}

\subsection*{What is proved}

The decomposition
$F^{\circ} = F^{\circ\circ} + \mathcal{M}$ is exact.
In the nontrivial cancellation regime $c_0 \ne 0$, the
character algebra proves $\alpha = 1$: the Mertens
divergence cancels algebraically, leaving $F^{\circ}$
as a pure non-trivial character sum. When $c_0 = 0$,
no cancellation is needed. Convergence
extends to the classical zero-free region for
$L$-functions modulo~$m$~\cite{davenport}. The
neutrality of the fixed mod-$3$ reflected class is
proved from the reflection identity.

\subsection*{What the computation shows}

At finite range ($348{,}488$ primes), the fitted
exponent $\alpha \approx 0.6$ reflects the
pre-asymptotic regime where character sum corrections
are comparable to the Mertens term. As the range
increases to $1{,}270{,}000$ primes, the fit drifts
toward $\alpha = 1$, confirming the theorem.

\subsection*{The arc}

The alignment limit $2/3$ from~\cite{paper1} is the
base-$3$ instance of the neutrality theorem: at base~$3$,
the class-$1$ bias is $-2/3$ and the class-$2$ mean is
$-1/3$ (neutral at $-1/2$ in the centered picture). The
universal structure is not the ratio $2/3$ but the
neutrality of the fixed reflected class, which follows
from the bilateral symmetry of the digit function. The collision
transform~\cite{paper16} provides the character
decomposition, and the antisymmetry
theorem~\cite{paper17} restricts to odd characters,
eliminating the Mertens term and producing perfect
cancellation.

% ============================================================
\begin{thebibliography}{99}

\bibitem{davenport}
H.~Davenport,
\emph{Multiplicative Number Theory},
3rd~ed., Springer, 2000.

\bibitem{hardy}
G.~H.~Hardy and M.~Riesz,
\emph{The General Theory of Dirichlet's Series},
Cambridge Tracts in Mathematics, 1915.

\bibitem{iwaniec}
H.~Iwaniec and E.~Kowalski,
\emph{Analytic Number Theory},
AMS Colloquium Publications, vol.~53, 2004.

\bibitem{montgomery}
H.~L.~Montgomery and R.~C.~Vaughan,
\emph{Multiplicative Number Theory I: Classical Theory},
Cambridge University Press, 2007.

\bibitem{paper1}
A.~S.~Petty,
Why the golden ratio selects the prime three,
research note, 2020.

\bibitem{paper16}
A.~S.~Petty,
The collision transform,
research note, 2025.

\bibitem{paper17}
A.~S.~Petty,
The collision transform in the critical strip,
research note, 2025.

\bibitem{nfield}
A.~S.~Petty,
\emph{nfield: A structural analysis engine for fractional
field invariants},
software.
\url{https://github.com/alexspetty/nfield}

\end{thebibliography}

\end{document}
